\scanpage{6}

cette condition est encore remplie si on stabilise en prenant des sommes directes de formes, et des formes induites sur sous-objets.

\inserted{\deleted{\unsure{Plus précis, pour $\varphi$ fixé, on a}} \unsure{Si|de $\sum_I$,} $I$ fini, les formes sur $\coprod M_i$ qui sont compatibles aux $\varphi_i$ sont celles qui sont compatibles \deleted{avec} à $\coprod_i\varphi_i$. Si $\varphi$ de Weil sur $M$, et $N\hookrightarrow M$, les formes sur $N$ comp. à $\varphi$ sont celles compatibles à $\varphi|N$}

7) \deleted{Appelons} Supposons $\mathcal M$ gradué, muni d'un ``objet de Tate'' $K(1)$ de degré $-2$. On appelle \uline{prépolarisation} \deleted{la donn} de $\mathcal M$ la donnée, \deleted{\unsure{d'une}} pour tt $M\in\operatorname{Ob}\mathcal M$ de degré $i$, d'une partie \uline{non vide} \deleted{\unsure{de}} $P(M)\subset\operatorname{Hom}(M\otimes M,K(-i))$, formé de formes $(-1)^i$-sym., \deleted{de Weil,} satisfaisant la condition

\[
\left\{\begin{aligned}
&\varphi\in P(M),\ \psi\in P(N)\quad(M,N\in\operatorname{Ob}\mathcal M^i),\ u:M\to N\quad\text{implique}\\
&\operatorname{Tr}u^{*(\varphi,\psi)}u>0\qquad\text{si }u\ne0.
\end{aligned}\right. \tag*{(P1)}
\]

Cela signifie que pour les $M\in\operatorname{Ob}\mathcal M^i$ ($i$ fixé) les $\varphi\in P(M)$ sont mutuellement compatibles.

\margin{\uline{NB} Si $\exists$ prépolarisation $P$, alors $\mathcal M$ est semi-simple}

Si $P$ et $P'$ sont des prépolarisations, alors $P\subset P'\overset{\text{déf}}{\Longleftrightarrow}\forall M$ de degré $i$, on a $P(M)\subset P'(M)$. Notion de prépolarisations maximales. Grâce à 4) on trouve

\uline{Proposition} Toute prépolarisation $P$ est contenue dans une unique prépolarisation maximale $\widehat P$. On a $\widehat P(M)=$ ensemble des formes de Weil sur $M$ compatibles aux formes de $P(M)$.

\uline{Cor.} a) Si $\varphi\in\widehat P(M)$, $\psi\in\widehat P(N)$, alors $\varphi\oplus\psi\in\widehat P(M\oplus N)$ ($M$ et $N$ de $=$ degré $i$)

b) Si $\varphi\in\widehat P(M)$ \deleted{\unsure{$\varphi\in\widehat P$}} $N\subset M$, alors $\varphi|N\in\widehat P(N)$

c) si $\varphi\in\widehat P(M)$, alors $\check\varphi(\unsure{-i|i})\in\widehat P(\check M\deleted{\unsure{)}}(-i))$

% archive sheet number: 5
